\documentclass[]{article}

\usepackage{amsmath}
\usepackage{multirow}
\usepackage{natbib}
\usepackage{graphicx} 


\begin{document}

\subsection{Simulated Dataset}
The analysis was performed on a dataset comprising 20 simulated damped harmonic oscillators. For each oscillator, the dataset provided time-series data for displacement ($x(t)$) and velocity ($v(t)$) over a 20-second interval. In addition to the time-series data, constant physical parameters were provided for each oscillator, including its mass ($m$), spring constant ($k$), and theoretical damping coefficient ($b$). The total measured energy at each time step, $E_{total}(t)$, was calculated from the instantaneous kinetic and potential energies and was corrupted by simulated measurement noise.

\subsection{Noise Characterization and Bias Deconvolution}
The core of our method is the analytical removal of a constant energy bias introduced by measurement noise. This bias arises because the energy calculation involves squared quantities, causing the variance of zero-mean noise to manifest as a positive energy offset. We assume that the physical energy of the oscillator has effectively decayed to zero in the late-time regime, allowing the signal in this window to be used for noise characterization.

For each oscillator, we isolated the displacement and velocity signals in the time interval $t > 15$ s. To ensure that our variance calculation was not contaminated by any residual physical decay, a local linear detrending was applied to both signals in this window. The empirical variances of the detrended displacement and velocity signals, denoted as $\sigma_x^2$ and $\sigma_v^2$ respectively, were then computed.

These variances were used to calculate the constant, noise-induced energy bias, $\Delta E_{\text{noise}}$, according to the formula for the total energy of a harmonic oscillator:
\begin{equation}
\Delta E_{\text{noise}} = \frac{1}{2} k \sigma_x^2 + \frac{1}{2} m \sigma_v^2
\end{equation}
A corrected energy trajectory, $E_{corrected}(t)$, was then generated by subtracting this bias term from the total measured energy:
\begin{equation}
E_{corrected}(t) = E_{total}(t) - \Delta E_{\text{noise}}
\end{equation}
To ensure physical realism, any resulting negative energy values, which can occur due to stochastic fluctuations, were clipped to zero, yielding a final non-negative energy trajectory.

\subsection{Validation and Evaluation Metrics}
The effectiveness of the deconvolution framework was validated by assessing its ability to recover the true physical damping rate from the corrected energy data. The theoretical damping rate, $\gamma_{theory}$, for each oscillator was calculated directly from its physical parameters as $\gamma_{theory} = b / (2m)$.

An observed damping rate, $\gamma_{obs}$, was determined for each oscillator by performing a non-linear least-squares fit of the corrected energy trajectory, $E_{corrected}(t)$, to the theoretical exponential decay model:
\begin{equation}
E(t) = E_0 \exp(-2\gamma_{obs} t)
\end{equation}
where $E_0$ is the initial energy, which was also a free parameter in the fit.

The performance of the method was quantified by analyzing the residuals between the observed and theoretical damping rates, $\Delta\gamma = \gamma_{obs} - \gamma_{theory}$. The primary evaluation metrics were the mean and standard deviation of these residuals, aggregated across all 20 oscillators in the dataset. A mean residual close to zero indicates the successful removal of systematic bias from the parameter estimation.

\end{document}
                